\documentclass[11pt,a4paper]{article}

\usepackage[margin=2.3cm]{geometry}
\usepackage{graphicx}
\usepackage{booktabs}
\usepackage{tabularx}
\usepackage{array}
\usepackage{multirow}
\usepackage{float}
\usepackage{caption}
\usepackage{amsmath}
\usepackage{enumitem}
\usepackage{xcolor}
\usepackage{listings}
\usepackage{tikz}
\usetikzlibrary{arrows.meta,positioning}
\usepackage{fancyhdr}
\usepackage{hyperref}
\hypersetup{colorlinks=true,linkcolor=blue!50!black,urlcolor=blue!50!black}

% ---- fonts & Persian support (polyglossia must come after hyperref) ----
\usepackage{fontspec}
\usepackage{polyglossia}
\setdefaultlanguage{english}
\setotherlanguage{persian}
\newfontfamily\persianfont[Script=Arabic,Scale=1.05]{Amiri}
\newfontfamily\persianfontsf[Script=Arabic,Scale=1.05]{Amiri}
\newfontfamily\persianfonttt[Script=Arabic,Scale=1.05]{Amiri}
\newcommand{\fa}[1]{\textpersian{#1}}

% ---- code listings ----
\definecolor{codebg}{gray}{0.96}
\definecolor{codekw}{RGB}{0,70,150}
\definecolor{codecm}{RGB}{90,120,90}
\definecolor{codestr}{RGB}{160,60,30}
\lstset{
  language=Python,
  basicstyle=\ttfamily\small,
  keywordstyle=\color{codekw}\bfseries,
  commentstyle=\color{codecm}\itshape,
  stringstyle=\color{codestr},
  backgroundcolor=\color{codebg},
  frame=single, rulecolor=\color{gray!40},
  breaklines=true, columns=fullflexible,
  showstringspaces=false, aboveskip=8pt, belowskip=8pt
}

\newcommand{\col}[1]{\texttt{#1}}
\setlist{itemsep=2pt,topsep=4pt}
\setlength{\emergencystretch}{2em}
\renewcommand{\arraystretch}{1.18}

\pagestyle{fancy}
\fancyhf{}
\lhead{Data Mining -- Mini-Project 1}
\rhead{Omid Basafa (403521078)}
\cfoot{\thepage}

% =====================================================================
\begin{document}

% ------------------------------ title page ---------------------------
\begin{titlepage}
  \centering
  \vspace*{1cm}
  \includegraphics[width=0.28\textwidth]{iust-logo.png}\par
  \vspace{0.8cm}
  {\Large Iran University of Science and Technology\par}
  \vspace{0.3cm}
  {\large School of Computer Engineering\par}
  \vspace{1.6cm}
  {\Huge\bfseries Data Mining\par}
  \vspace{0.6cm}
  {\LARGE Mini-Project 1: Data Autopsy\par}
  \vspace{0.3cm}
  {\large Exploring, auditing and cleaning \texttt{iran\_poi\_raw.csv}\par}
  \vfill
  \begin{tabular}{rl}
    \textbf{Student name:}   & Omid Basafa \\[4pt]
    \textbf{Student number:} & 403521078 \\
  \end{tabular}
  \vfill
  {\large \today\par}
\end{titlepage}

\tableofcontents
\newpage

% =====================================================================
\setcounter{section}{-1}   % Introduction = 0, so section numbers match the step numbers
\section{Introduction}

This report describes the analysis of \col{iran\_poi\_raw.csv}, a raw, intentionally
uncleaned dataset of \textbf{297{,}152 real points of interest (POIs) in Iran}, such as
restaurants, mosques, schools, banks, parks and doctors' offices.
The data was extracted from OpenStreetMap (OSM) through the Geofabrik extract.
It was entered by volunteer OSM mappers from 2007 until today, with no central quality
control. The file is about 48\,MB and has 28 columns.

All work was done in a single Jupyter notebook (\col{main.ipynb}) with
\textbf{Python 3}, \textbf{pandas}, \textbf{NumPy}, \textbf{Matplotlib} and the standard
modules \col{re} (regular expressions) and \col{json}.
For each step this report gives the method, the code idea, the output, and an interpretation
of the result.

\subsection{The dataset}

\begin{table}[H]
\centering
\caption{Columns of \col{iran\_poi\_raw.csv}.}
\begin{tabularx}{\textwidth}{>{\raggedright\arraybackslash}p{4.3cm}X}
\toprule
\textbf{Column(s)} & \textbf{Meaning} \\
\midrule
\col{osm\_type}, \col{osm\_id} & Identify the OSM object (node or way, and its id). \\
\col{version} & Number of times the object has been edited (the maximum in the file is 50). \\
\col{edit\_date} & Date of the \emph{last} change. \\
\col{lat}, \col{lon} & Geographic position. \\
\col{main\_key}, \col{main\_value} & Type of the place, e.g.\ \col{amenity} (category) / \col{restaurant} (sub-category). \\
\col{name}, \col{name:fa}, \col{name:en} & Names, exactly as typed by the mapper. \\
\col{addr:city}, \col{addr:street}, \col{addr:housenumber} & Address, entered without any fixed rule. \\
\col{phone}, \col{website}, \col{opening\_hours} & Free-text contact and opening-hours fields. \\
\col{brand}, \col{operator}, \col{cuisine}, \col{wheelchair}, \col{building}, \col{religion} & Optional descriptive tags. \\
\col{capacity} & In theory the capacity of the place as a number. \\
\col{n\_tags} & Total number of OSM tags of the record. \\
\col{nearest\_hub}, \col{hub\_km} & Added by the course staff (see the note below). \\
\col{other\_tags} & All remaining OSM tags, stored as a JSON object. \\
\bottomrule
\end{tabularx}
\end{table}

\noindent\fbox{\parbox{0.97\textwidth}{\textbf{Important note on \col{nearest\_hub}.}
\col{nearest\_hub} assigns each record to the \emph{nearest approximate provincial center}
using \emph{only straight-line geographic distance}, and \col{hub\_km} is that distance in km.
This is \textbf{not an official administrative division}: a record labelled ``Tehran'' is not
necessarily inside Tehran province. Every result in this report that uses these columns
(Steps 6.2--6.4) should be read with this in mind.}}

\vspace{0.4cm}
A few input rows (selected columns) as loaded by \col{pd.read\_csv}:

\begin{table}[H]
\centering\footnotesize
\setlength{\tabcolsep}{4pt}
\caption{Sample input records (\col{df.head()}, selected columns).}
\begin{tabular}{llrlllll}
\toprule
\col{osm\_type} & \col{osm\_id} & \col{version} & \col{edit\_date} & \col{main\_key} & \col{main\_value} & \col{name} & \col{nearest\_hub} \\
\midrule
node & 5        & 15 & 2026-08-01 & amenity & bureau\_de\_change & \fa{صرافی MP} & Tehran \\
node & 25990999 & 5  & 2019-08-27 & amenity & toilets & \fa{دستشویی} & Tehran \\
node & 26016253 & 10 & 2019-07-18 & amenity & pharmacy & \fa{داروخانه شبانه‌روزی قانون} & Tehran \\
node & 26206700 & 9  & 2019-01-26 & amenity & cinema & \fa{سینما عصر جدید} & Tehran \\
\bottomrule
\end{tabular}
\end{table}

\subsection{Overall workflow}

\begin{figure}[H]
\centering
\begin{tikzpicture}[
  box/.style={draw, rounded corners, fill=blue!6, align=center, minimum height=1.25cm,
              text width=2.25cm, font=\small},
  arr/.style={-{Stealth[length=2.5mm]}, thick}]
  \node[box, fill=gray!12] (raw) {Raw CSV\\297{,}152 rows\\28 columns};
  \node[box, right=0.45cm of raw] (s2) {Step 2\\types \&\\basic stats};
  \node[box, right=0.45cm of s2] (s3) {Step 3\\quality audit};
  \node[box, right=0.45cm of s3] (s4) {Step 4\\cleaning};
  \node[box, below=1.1cm of s4] (s5) {Step 5\\Jaccard \&\\duplicates};
  \node[box, left=0.45cm of s5] (s6) {Step 6\\statistics,\\plots, OLAP};
  \node[box, left=0.45cm of s6, fill=green!8] (rep) {Results \&\\report};
  \node[box, left=0.45cm of rep, fill=yellow!10] (s1) {Step 1\\task\\identification};
  \draw[arr] (raw) -- (s2);
  \draw[arr] (s2) -- (s3);
  \draw[arr] (s3) -- (s4);
  \draw[arr] (s4) -- (s5);
  \draw[arr] (s5) -- (s6);
  \draw[arr] (s6) -- (rep);
  \draw[arr, dashed] (s3.south) to[bend right=25] node[midway, right, font=\scriptsize]{duplicate candidates} (s5.north west);
\end{tikzpicture}
\caption{Overall workflow. Step 1 is theoretical. The duplicate candidates found in
Step 3.5 are saved to \col{candids\_duplicate.csv} and reused in Step 5.}
\end{figure}

The file is loaded with \col{phone} and \col{capacity} forced to \emph{string}:

\begin{lstlisting}
df = pd.read_csv("iran_poi_raw.csv", dtype={"capacity": str, "phone": str})
\end{lstlisting}

\noindent Without this, pandas would try to infer a numeric type: leading zeros
and the \col{+} sign of phone numbers would be lost, and the mixed \col{capacity} column
would be guessed into some type. We want to inspect the raw values first and clean them
ourselves (Step 4).

% =====================================================================
\section{Step 1: Identifying the Data-Mining Tasks}

\begin{table}[H]
\centering
\caption{Data-mining task for each of the three questions.}
\begin{tabularx}{\textwidth}{>{\raggedright\arraybackslash}p{4.4cm}>{\raggedright\arraybackslash}p{2.5cm}X}
\toprule
\textbf{Question} & \textbf{Task} & \textbf{Reason} \\
\midrule
Can the type of a place (\col{main\_value}) be guessed from its name alone?
& \textbf{Classification}
& The target \col{main\_value} is a \emph{categorical label} that is already known for most
records, so a model can be trained with supervision on (name $\rightarrow$ type) examples
and then predict the type of a new name. \\
\midrule
Which geographic regions have a higher density of recreational places?
& \textbf{Clustering}
& There is no label to predict. We want to discover dense groups of
\col{leisure}/\col{tourism} points in (\col{lat}, \col{lon}) space without supervision,
e.g.\ with a density-based algorithm such as DBSCAN. \\
\midrule
Which records are duplicate versions of another place?
& \textbf{Anomaly detection}
& Normally each real place appears once. A record that almost coincides with another one
(same location, same type, similar name) is an \emph{abnormal} record that should be flagged.
It can also be seen as a similarity-based grouping problem, which is how Step 5 solves it. \\
\bottomrule
\end{tabularx}
\end{table}

% =====================================================================
\section{Step 2: Attribute Types and Basic Statistics}

The data was first inspected with \col{df.shape}, \col{df.head()}, \col{df.info()},
\col{df.dtypes} and \col{df.describe()}.

\subsection{Attribute types}

\begin{table}[H]
\centering
\caption{Attribute type of the requested columns.}
\begin{tabularx}{\textwidth}{>{\raggedright\arraybackslash}p{2.6cm}>{\raggedright\arraybackslash}p{2cm}X}
\toprule
\textbf{Column} & \textbf{Type} & \textbf{Reason} \\
\midrule
\col{main\_key} & Nominal & Category labels (amenity, shop, \dots) with no natural order; only $=$ / $\neq$ make sense. \\
\col{version} & Ratio & A count of edits. It has a true zero and ratios make sense (10 edits is twice as many as 5). \\
\col{edit\_date} & Interval & Differences between dates are meaningful (``3 years later''), but the calendar zero is arbitrary, so ratios of dates are meaningless. \\
\col{lat} / \col{lon} & Interval & Differences in degrees are meaningful, but the zero point (equator / Greenwich) is a convention and values can be negative, so ratios are meaningless. \\
\col{capacity} & Ratio & A number of people or vehicles. 0 really means ``no capacity'' and 100 is twice 50. \\
\col{hub\_km} & Ratio & A distance in km, with a true zero (being at the hub) and meaningful ratios. \\
\bottomrule
\end{tabularx}
\end{table}

\subsection{Basic statistics}

\textbf{Total number of records: 297{,}152} (28 columns).
The share of each of the seven \col{main\_key} categories was computed with
\col{value\_counts()} divided by the number of rows:

\begin{table}[H]
\centering
\caption{Share of each \col{main\_key} category.}
\begin{tabular}{lrr}
\toprule
\col{main\_key} & \textbf{Count} & \textbf{Percentage} \\
\midrule
amenity    & 138{,}493 & 46.61\% \\
shop       & 72{,}122  & 24.27\% \\
leisure    & 42{,}320  & 14.24\% \\
office     & 25{,}972  & 8.74\%  \\
tourism    & 12{,}220  & 4.11\%  \\
craft      & 3{,}911   & 1.32\%  \\
healthcare & 2{,}114   & 0.71\%  \\
\midrule
\textbf{Total} & \textbf{297{,}152} & \textbf{100\%} \\
\bottomrule
\end{tabular}
\end{table}

The percentage of missing values per column was computed with:

\begin{lstlisting}
missing = (df.isna().mean() * 100).round(2).sort_values(ascending=True)
\end{lstlisting}

\begin{table}[H]
\centering
\caption{Percentage of missing values in each column.}
\begin{tabular}{lr@{\hspace{1.6cm}}lr}
\toprule
\textbf{Column} & \textbf{Missing \%} & \textbf{Column} & \textbf{Missing \%} \\
\midrule
\col{capacity}         & 99.71 & \col{opening\_hours} & 88.34 \\
\col{wheelchair}       & 99.23 & \col{addr:city}      & 84.81 \\
\col{cuisine}          & 98.69 & \col{name:en}        & 84.62 \\
\col{name:fa}          & 97.55 & \col{phone}          & 84.35 \\
\col{addr:housenumber} & 96.95 & \col{addr:street}    & 82.02 \\
\col{website}          & 96.12 & \col{other\_tags}    & 65.10 \\
\col{brand}            & 95.35 & \col{name}           & 24.76 \\
\col{religion}         & 94.14 & \multicolumn{2}{l}{\textit{The other 11 columns: 0.00}} \\
\col{operator}         & 92.00 & \multicolumn{2}{l}{\small(\col{osm\_type}, \col{osm\_id}, \col{version}, \col{edit\_date},} \\
\col{building}         & 88.86 & \multicolumn{2}{l}{\small\phantom{(}\col{lat}, \col{lon}, \col{main\_key}, \col{main\_value},} \\
                       &       & \multicolumn{2}{l}{\small\phantom{(}\col{n\_tags}, \col{nearest\_hub}, \col{hub\_km})} \\
\bottomrule
\end{tabular}
\end{table}

\begin{figure}[H]
\centering
\includegraphics[width=0.62\textwidth]{figures/missing.png}
\caption{Missing percentage per column.}
\end{figure}

Only the fields that are always produced by the extraction are complete. All the
descriptive tags are mostly empty, and even \col{name} is missing in about a quarter of
the records.

\subsection{Why are \col{capacity} and \col{cuisine} more than 95\% empty?}

To check whether the empty values are random, we looked at the overall fill rate of
\col{cuisine} (only \textbf{1.31\%} of all records) and at which kinds of places actually
have a \col{capacity}:

\begin{table}[H]
\centering
\caption{\col{main\_value} of the 854 records that have a \col{capacity} (top rows).}
\begin{tabular}{lr@{\hspace{1.2cm}}lr}
\toprule
\col{main\_value} & \textbf{Count} & \col{main\_value} & \textbf{Count} \\
\midrule
parking        & 524 & shelter            & 22 \\
restaurant     & 87  & fast\_food         & 17 \\
parking\_space & 70  & camp\_site         & 12 \\
cafe           & 33  & school             & 11 \\
taxi           & 27  & \textit{others}    & 51 \\
\bottomrule
\end{tabular}
\end{table}

A high missing rate here is \textbf{not necessarily a data-collection error}:
\begin{itemize}
  \item \textbf{Not applicable (structurally missing).} \col{cuisine} only makes sense for
  food places (restaurant, cafe, fast\_food). A mosque, a bank or a park has no cuisine, so an
  empty value is the \emph{correct} value there. In the same way \col{capacity} is mostly
  used for parking (524 of 854), \col{religion} only for places of worship, and \col{brand}
  only for chain stores.
  \item \textbf{Optional tags in a volunteer project.} In OSM only the main tag is
  essential. Every other tag is optional, and many mappers only add a point and a name.
  Details stay empty even when they would apply.
\end{itemize}
So the values are \emph{missing not at random}: whether a value is present depends on the
type of the place and on how much effort the mapper put in, not on a failure of the
collection system.

% =====================================================================
\section{Step 3: Data Quality}

\subsection{How messy is \col{addr:city}?}

Only 15.19\% of the records have \col{addr:city}, yet the column has
\textbf{2{,}525 unique values}, far more than the number of cities in Iran.
All values containing a spelling of Tehran were found with a case-insensitive regex
(\col{tehran|teheran|\fa{تهران}|\fa{طهران}}):

\begin{table}[H]
\centering
\caption{Different written forms of the city of Tehran in \col{addr:city}.}
\begin{tabular}{lr@{\hspace{1.2cm}}lr}
\toprule
\textbf{Value} & \textbf{Count} & \textbf{Value} & \textbf{Count} \\
\midrule
\fa{تهران}     & 6{,}110 & \col{tehran-pakdasht-hesaramir} & 12 \\
\col{Tehran}     & 1{,}018 & \col{TEHRAN}                   & 8 \\
\fa{شهر تهران} & 932     & \fa{منطقه ۱۴ شهر تهران}        & 6 \\
\col{tehran}     & 408     & \fa{استان تهران}               & 5 \\
\col{Teheran}    & 13      & \fa{تهران تجریش}               & 4 \\
\bottomrule
\end{tabular}
\end{table}

At least six different spellings refer to exactly the same city:
\fa{تهران}, \col{Tehran}, \fa{شهر تهران}, \col{tehran}, \col{Teheran} and \col{TEHRAN}.
There are also district or province forms, combined forms, and even complete addresses
pasted into the city field (e.g.\ \fa{تهران، باغ سپهسالار، پاساژ الیزه، ...}).
The substring search also matched \fa{تهراندشت} and \fa{تهرانسر}, which are
\emph{different} places. This shows that simple substring matching is not enough for cleaning.

\subsection{Arabic \fa{ك}/\fa{ي} instead of Persian \fa{ک}/\fa{ی}}

The two forms look almost identical, but they are different Unicode code points
(U+0643/U+064A versus U+06A9/U+06CC), so a computer treats \fa{بانك} and \fa{بانک}
as different strings. Counting with \col{str.contains} on the \col{name} column gives:

\begin{center}
\begin{tabular}{lr}
\toprule
Names containing Arabic \fa{ك} & \textbf{765} \\
Names containing Arabic \fa{ي} & \textbf{1{,}486} \\
\bottomrule
\end{tabular}
\end{center}

\textbf{Why is this noise specific to user-generated data?}
Volunteers type on different devices and keyboard layouts. The Arabic layout of
Windows, older phones and some apps produces \fa{ي}/\fa{ك}, while the standard
Persian layout produces \fa{ی}/\fa{ک}. Because both look the same on screen, the mapper
does not notice, and OSM does not validate free text. An organization that produces data in a
fixed format uses one input system and applies validation and normalization rules, so this kind of
mixing does not appear. As a result, searching, grouping and string matching treat
\fa{بانك ملي} and \fa{بانک ملی} as different values.

\subsection{Invalid values in \col{capacity}}

Each of the 854 non-null \col{capacity} strings was put into a category with a function
\col{valid\_capacity} built from ordered regular-expression checks (the first match wins):
a plain integer, an integer in Persian digits, a negative number, digits followed by a unit
(``\fa{نفر}'' = person, ``\fa{متر}'' = metre), digits with other characters, and no digits at all.

\begin{table}[H]
\centering
\caption{Categories of \col{capacity} values.}
\begin{tabular}{lr}
\toprule
\textbf{Category} & \textbf{Count} \\
\midrule
valid integer (ok)       & 842 \\
Persian digits           & 4 \\
text instead of number   & 3 \\
negative number          & 2 \\
number + unit            & 2 \\
number + extra letters   & 1 \\
\bottomrule
\end{tabular}
\end{table}

\begin{table}[H]
\centering
\caption{Invalid \col{capacity} values and their probable source.}
\begin{tabularx}{\textwidth}{>{\raggedright\arraybackslash}p{2.6cm}>{\raggedright\arraybackslash}p{4.1cm}X}
\toprule
\textbf{Category} & \textbf{Examples (place)} & \textbf{Probable source} \\
\midrule
Persian digits & \fa{۵۰} (cafe), \fa{۲۰}, \fa{۶۰} (parking), \fa{۱۰} (cafe)
& Typed with a Persian keyboard. The value is correct but not machine-readable as a number. \\
Negative number & $-1$ (taxi), $-50$ (restaurant)
& $-1$ is a common placeholder for ``unknown'', and $-50$ is probably a typo. A capacity cannot be negative. \\
Number + unit & \fa{۵ نفر} (5 persons, taxi), \fa{۸۰۰متر مربع} (800\,m$^2$, parking)
& The user wrote a phrase instead of a number. The second one is even the \emph{wrong
quantity}: an area instead of a number of people or cars. \\
Number + extra letters & \fa{ضض}300 (parking space)
& A keyboard typo: the layout was not switched before typing (\fa{ض} is on the same key as \texttt{q}). \\
Text instead of number & \fa{مشهد} (a city), \fa{ترمینال}, \fa{مخزن (تانکر) آب حسینی} (the place's own name)
& Content typed into the wrong field of the editor: the city or the name was put into \col{capacity}. \\
\bottomrule
\end{tabularx}
\end{table}

Some syntactically valid integers are also \emph{logically} suspicious: 5 records have a
capacity of \textbf{0}, and some values are implausible, for example an \emph{internet-taxi office}
with 10{,}000 and a \emph{restaurant} with 5{,}000. Others, such as a 50{,}000-car parking for
Arbaeen or a 15{,}000-seat stadium, are plausible.

\subsection{Phone formats}

Each phone value was assigned to a format by \col{format\_phone}.
A helper \col{nationals} reduces a number to its \emph{national significant number}: it keeps
only the digits, removes the country code (\col{0098} or \col{+98}) and strips the leading trunk
\col{0}. A complete Iranian number then has 10 digits (e.g.\ \col{2126119186}).

\begin{lstlisting}
def nationals(p):
    d = re.sub(r"\D", "", p)              # keep digits only
    if d.startswith("0098"):
        d = d[4:]
    elif d.startswith("98") and p.strip().startswith("+"):
        d = d[2:]
    return d.lstrip("0")                  # drop trunk prefix 0
\end{lstlisting}

\begin{table}[H]
\centering
\caption{Phone formats among the 46{,}490 non-null phones, with real samples.}
\begin{tabular}{lrl}
\toprule
\textbf{Format} & \textbf{Count} & \textbf{Real samples} \\
\midrule
Starts with 0, no \col{+98}            & 15{,}168 & \col{09111153844}, \col{08733779673} \\
\col{+98} without spaces               & 14{,}537 & \col{+983132205774}, \col{+983133375777} \\
\col{+98} with spaces                  & 11{,}537 & \col{+98 21 4423 8400}, \col{+98 21 4425 3980} \\
Persian / Arabic digits                & 2{,}735  & \fa{۰۲۱۴۹۲۸۶۰۰۰} \\
Several numbers in one field           & 1{,}113  & \col{+98 21 2611 9186;+98 21 2611 9187} \\
Neither 0 nor \col{+98}                & 554      & \col{9334440337}, \col{+5138412829} \\
Incomplete ($<8$ national digits)      & 492      & \col{+98 21 64451}, \col{+98 21 1836} \\
With \col{-} or \col{()}               & 339      & \col{+98-511-7654476-9}, \col{+98-511-8545061-3} \\
Not a phone number                     & 15       & \fa{رضا اسدی} (a person's name), \fa{کلینیک پوست و مو باران} \\
\bottomrule
\end{tabular}
\end{table}

The same information appears in at least nine formats. The variations are:
\begin{itemize}
  \item \textbf{digits:} Persian or Latin;
  \item \textbf{prefix:} \col{+98}, \col{0}, or none (\col{+5138412829} has lost the \col{98});
  \item \textbf{separators:} spaces, hyphens or parentheses;
  \item \textbf{completeness:} numbers cut off after the area code;
  \item \textbf{other content:} several numbers in one field, or plain text.
\end{itemize}
\col{+98-511-...} also uses the \emph{old} 3-digit Mashhad area code.

\subsection{Candidate duplicate records}

\col{lat} and \col{lon} were rounded to 4 decimals (about 11\,m in latitude and 9--10\,m in
longitude at Iran's latitudes). Records with the same rounded location \emph{and} the same
\col{main\_value} were grouped. Every group of size $>1$ is a set of duplicate candidates,
and a self-join turns the groups into pairs:

\begin{lstlisting}
df["lat_round"] = df["lat"].round(4)
df["lon_round"] = df["lon"].round(4)
keys = ["lat_round", "lon_round", "main_value"]
group_size = df.groupby(keys)["osm_id"].transform("count")
candids_duplicate = df[group_size > 1]

d = candids_duplicate.reset_index()
pairs = d.merge(d, on=keys, suffixes=("_1", "_2"))
pairs = pairs[pairs["index_1"] < pairs["index_2"]]   # no (A,A), no (B,A)
pairs.to_csv("candids_duplicate.csv", index=False)   # kept for Step 5
\end{lstlisting}

\begin{center}
\begin{tabular}{lr}
\toprule
Records in a duplicate group & \textbf{7{,}076} \\
Number of groups & \textbf{3{,}245} \\
Candidate pairs & \textbf{4{,}747} \\
Pairs with exactly the same name & \textbf{460} \\
\bottomrule
\end{tabular}
\end{center}

\begin{table}[H]
\centering
\caption{Real examples of the same place mapped twice.}
\begin{tabular}{llll}
\toprule
\textbf{Name} & \col{main\_value} & \textbf{Record 1} & \textbf{Record 2} \\
\midrule
\fa{کلانتری ۱۲ میدان جمهوری} & police      & node 455969203  & way 929870598 \\
\fa{بانک تجارت}              & bank        & node 2309413230 & way 1419976030 \\
\fa{اداره پست مرکزی}         & post\_office & node 3303523321 & way 491120128 \\
\fa{داروخانه دکتر بهشتیان}   & pharmacy    & node 3343366958 & way 327449401 \\
\fa{هتل نویذر}               & hotel       & node 3300670011 & node 5198144121 \\
\bottomrule
\end{tabular}
\end{table}

A common pattern is \textbf{node + way}: one mapper put a point (node) for the place,
and later another mapper drew the building outline (way) and tagged it again without
removing the original point. The hotel example shows that two plain nodes can also be
duplicates. The full candidate list was saved to \col{candids\_duplicate.csv} for Step 5.

% =====================================================================
\section{Step 4: Pre-processing and Cleaning}

\subsection{Normalizing \col{addr:city}}

The function \col{normalize} applies, in order:
\begin{enumerate}
  \item collapse repeated, leading and trailing whitespace;
  \item replace Arabic \fa{ي}/\fa{ك} with Persian \fa{ی}/\fa{ک};
  \item lower-case Latin letters (\col{Tehran}, \col{TEHRAN} $\rightarrow$ \col{tehran});
  \item remove the prefix ``\fa{شهر}'' (``city'');
  \item map the Latin name of \textbf{25 frequent cities} (more than the 10 required) to its
  Persian form with a dictionary, e.g.\ \col{tehran}/\col{teheran} $\rightarrow$ \fa{تهران},
  \col{isfahan}/\col{esfahan} $\rightarrow$ \fa{اصفهان}, \col{mashhad} $\rightarrow$ \fa{مشهد},
  \col{shiraz} $\rightarrow$ \fa{شیراز}, \col{tabriz}, \col{karaj}, \col{ahvaz}/\col{ahwaz},
  \col{qom}, \col{rasht}, \col{kermanshah}, \dots
\end{enumerate}

\begin{table}[H]
\centering
\caption{Effect of normalization.}
\begin{tabular}{lr}
\toprule
Unique values \textbf{before} & 2{,}525 \\
Unique values \textbf{after}  & 2{,}281 \\
Reduction & 244 (9.7\%) \\
\bottomrule
\end{tabular}
\hspace{1cm}
\begin{tabular}{lr}
\toprule
\textbf{Original} $\rightarrow$ \fa{تهران} & \textbf{Count} \\
\midrule
\fa{تهران} & 6{,}110 \\
\col{Tehran} & 1{,}018 \\
\fa{شهر تهران} & 932 \\
\col{tehran} & 408 \\
\col{Teheran} & 13 \\
\col{TEHRAN} & 8 \\
\bottomrule
\end{tabular}
\end{table}

All six spellings of Tehran now become a single value (8{,}489 records).
The remaining variety (district names, full addresses, villages) cannot be fixed with simple
rules and would need fuzzy matching or a list of official place names.

\subsection{Converting \col{capacity} to an integer}

The function \col{parse\_cap}:
\begin{enumerate}
  \item converts Persian (\fa{۰}--\fa{۹}) and Arabic-Indic digits to Latin digits with a
  single \col{str.maketrans} translation table;
  \item removes the unit ``\fa{نفر}'' (persons), since a number of persons is still a valid capacity;
  \item accepts the value only if what remains is a plain non-negative integer, checked with
  the regex \col{[0-9]+} and \col{re.fullmatch}; otherwise it returns \col{NaN}.
\end{enumerate}
The result is stored as pandas' nullable integer type \col{Int64}, so integers and missing
values can be in the same column.

\begin{table}[H]
\centering
\caption{Result of the \col{capacity} conversion.}
\begin{tabular}{lr}
\toprule
Non-null values before & 854 \\
Converted to integer & \textbf{847} \\
Set to \col{NaN} (cannot be converted logically) & \textbf{7} \\
\bottomrule
\end{tabular}
\end{table}

The 847 values are the 842 valid integers, the 4 Persian-digit values, and
\fa{۵ نفر}$\,\rightarrow 5$. The 7 \col{NaN}s are:
\begin{itemize}
  \item the two negative numbers ($-1$, $-50$);
  \item \fa{۸۰۰متر مربع}, which has the wrong unit (an area, not a person count);
  \item \fa{ضض}300, where the intended number is uncertain;
  \item the three text values (\fa{مشهد}, \fa{ترمینال}, \fa{مخزن (تانکر) آب حسینی}).
\end{itemize}

\subsection{Unifying \col{phone}}

All phone digits were converted to Latin digits with the same translation table.
The rule for a \textbf{``probably invalid'' number} is: after the conversion and after
removing the country code and the leading 0 (\col{nationals}), a complete Iranian number has
10 digits, so a number with \textbf{fewer than 8 digits} cannot be complete. A field may hold
several numbers, so it is split on \col{;}, \col{,} and \col{/}, and the field is flagged if
\emph{any} part is too short:

\begin{lstlisting}
def is_invalid(p):
    for part in re.split(r"[;,/]", p):
        if len(nationals(part)) < 8:
            return True
    return False
\end{lstlisting}

\begin{table}[H]
\centering
\caption{Result of phone unification.}
\begin{tabular}{lr}
\toprule
Phones with Persian/Arabic digits, before & 2{,}742 \\
Phones with Persian/Arabic digits, after  & 0 \\
\textbf{Probably invalid phones} & \textbf{550 of 46{,}490 (1.18\%)} \\
\bottomrule
\end{tabular}
\end{table}

Typical invalid numbers were cut off after the area code, such as \col{+98 21 5418 5},
\col{+98 13 3611} and \col{+982123540}. The 15 text-only values are also flagged, because they have no digits.

% =====================================================================
\section{Step 5: Similarity and Distance}

\subsection{Jaccard similarity of \col{other\_tags} keys}

The Jaccard similarity of two sets is
\[
J(A,B) = \frac{|A\cap B|}{|A\cup B|}.
\]
As required, it was implemented by hand with Python set operations (\col{\&} for
intersection, \col{|} for union), without any library function. When both sets are empty,
the union is empty and the similarity is undefined, so \col{NaN} is returned instead of
dividing by zero:

\begin{lstlisting}
def jaccard(a, b):
    a, b = set(a), set(b)
    union = a | b
    if len(union) == 0:          # both empty -> undefined
        return np.nan
    return len(a & b) / len(union)

def tag_keys(j):                 # set of keys of the other_tags JSON
    if pd.isna(j):
        return set()
    return set(json.loads(j).keys())
\end{lstlisting}

A sanity check gives $J(\{a,b,c\},\{b,c,d\}) = 2/4 = 0.5$.
From the 4{,}747 candidate pairs of Step 3.5, all pairs in which \textbf{both} records have a
non-empty \col{name} were selected: \textbf{2{,}396 pairs}, far more than the 20 required.
For each pair \col{tag\_sim} was computed as the Jaccard similarity of the \col{other\_tags}
key sets. For 772 pairs neither record has any \col{other\_tags}, so \col{tag\_sim} is \col{NaN}.

\subsection{Jaccard similarity of name words}

For the same pairs, each name was turned into a \emph{set of words} by splitting on
whitespace (e.g.\ \fa{بانک ملی} $\rightarrow$ \{\fa{بانک}, \fa{ملی}\}).
Before splitting, Arabic \fa{ي}/\fa{ك} were replaced by Persian letters and Latin text was
lower-cased, so spelling noise does not hide real matches. \col{name\_sim} is the Jaccard
similarity of the two word sets.

\begin{table}[H]
\centering
\caption{A sample of the compared pairs.}
\begin{tabular}{llcc}
\toprule
\col{name\_1} & \col{name\_2} & \col{tag\_sim} & \col{name\_sim} \\
\midrule
\fa{بیمارستان شهدای کردیان} & \fa{بیمارستان شهدای کردیان} & 1.00 & 1.00 \\
\fa{بانک کشاورزی} & \fa{بانک کشاورزی} & 0.00 & 1.00 \\
\fa{مسجد بلال} & \fa{مسجد بلال} & 0.00 & 1.00 \\
\fa{نی نی پوش} & \fa{پوشاک ایزدی} & 1.00 & 0.00 \\
\fa{دکتر محمدرضا اهتمام} & \fa{دکتر لیلا توکلی} & 1.00 & 0.20 \\
\fa{بانک تجارت} & \fa{بانک سپه} & 0.50 & 0.33 \\
\fa{بونیک مارکت شعبه 1 اندیشه} & \fa{بونیک مارکت شعبه 2 اندیشه} & 1.00 & 0.67 \\
\fa{پارک شهدا} & \fa{پارک مردم)(شهدا)} & NaN & 0.33 \\
\bottomrule
\end{tabular}
\end{table}

\paragraph{Do pairs with high tag similarity also have high name similarity?}
\textbf{No.} The two measures are practically unrelated:

\begin{table}[H]
\centering
\caption{Relation between \col{tag\_sim} and \col{name\_sim}.}
\begin{tabular}{lr}
\toprule
Pearson correlation of \col{tag\_sim} and \col{name\_sim} & $-0.063$ \\
Mean \col{name\_sim} when \col{tag\_sim} $\geq 0.8$ & 0.30 \\
Mean \col{name\_sim} when \col{tag\_sim} $< 0.8$ & 0.32 \\
\bottomrule
\end{tabular}
\hspace{0.6cm}
\begin{tabular}{lcc}
\toprule
 & name $\leq 0.5$ & name $> 0.5$ \\
\midrule
tag $< 0.5$    & 502 & 167 \\
tag $\geq 0.5$ & 788 & 167 \\
\bottomrule
\end{tabular}
\end{table}

\begin{figure}[H]
\centering
\includegraphics[width=0.55\textwidth]{figures/sim_scatter.png}
\caption{Tag similarity versus name similarity (one point per pair). There is no diagonal
trend. The points form vertical stripes at simple fractions (0, 1/3, 1/2, 2/3, 1) because the
tag key sets are very small.}
\end{figure}

\paragraph{Pairs on which the two measures disagree.}
\begin{enumerate}
  \item \textbf{Same name, tag similarity 0 (93 pairs).} Example: \fa{اداره پست مرکزی}
  (node 3303523321 / way 491120128). The node has no \col{other\_tags}, and the way has
  \col{\{"building:levels": "5"\}}. Another example is \fa{شبکه تلویزیونی پرس تی‌وی}, whose way has
  \col{building:levels} and \col{operator:wikidata}.
  These are node + way duplicates: the way (the building polygon) carries
  \emph{geometry/building} tags that the point does not have. The tag keys describe
  \emph{how} the object was mapped, not \emph{what} it is, so with one empty set the
  similarity is 0 even though this is clearly the same place.
  \item \textbf{Identical tag keys, different names (401 pairs).} Example:
  \fa{بانک اقتصاد نوین} versus \fa{بانک ایران زمین}. Both have the keys
  \col{brand:en}, \col{brand:wikidata}, \dots, but with different \emph{values}
  (EN Bank versus Iran Zamin Bank). Two clinics both tagged \col{healthcare=clinic}
  and two fast-food shops both tagged \col{delivery}, \col{delivery:partner} behave the same way.
  Places of the same type get the \emph{same standard set of keys}, often from the same
  editor preset, and only keys are compared. These are \textbf{different businesses} that happen
  to be next to each other, which is common in shopping centres and medical buildings.
\end{enumerate}

\textbf{Conclusion.} The Jaccard similarity of tag \emph{keys} measures how similarly two
records were \emph{tagged}, not whether they are the same place. The Jaccard similarity of name
words is a much more reliable signal for duplicate detection.

\subsection{Duplicate threshold}

To choose a threshold, the name similarities were divided into bands, and the borderline
band was checked by hand:

\begin{table}[H]
\centering
\caption{Distribution of \col{name\_sim} over the 2{,}396 named pairs.}
\begin{tabular}{lr}
\toprule
\col{name\_sim} band & \textbf{Pairs} \\
\midrule
$< 0.5$            & 1{,}749 \\
0.5 -- 0.65        & 115 \\
0.66 -- 0.99       & 66 \\
1 (same words)     & 466 \\
\bottomrule
\end{tabular}
\end{table}

\begin{table}[H]
\centering
\caption{Examples from the borderline band 0.5--0.65.}
\begin{tabular}{llcc}
\toprule
\col{name\_1} & \col{name\_2} & \col{tag\_sim} & \col{name\_sim} \\
\midrule
\fa{طلا و جواهر اکسید} & \fa{طلا و جواهر تک} & 1.0 & 0.6 \\
\fa{مطب دکتر اصغری} & \fa{مطب دکتر سخاوت‌فر} & 1.0 & 0.5 \\
\fa{آش سرای برق} & \fa{کباب سرای برق} & 0.0 & 0.5 \\
\fa{اغذیه ی پونک} & \fa{اغذیه ی آوین} & NaN & 0.5 \\
\fa{آزمایشگاه بابک} & \fa{آزمایشگاه تشخیص پزشکی بابک} & 0.0 & 0.5 \\
\fa{مسجد نقویه} & \fa{مسجد و حسینیه نقویه} & NaN & 0.5 \\
\bottomrule
\end{tabular}
\end{table}

This band is mostly \textbf{false positives}. Different shops share generic words such as
\fa{طلا و جواهر} (jewellery) or \fa{مطب دکتر} (doctor's office), while the distinctive
part of the name differs. Only a few pairs are true duplicates (the laboratory and mosque
examples). From about two-thirds of shared words upward, the pairs are mostly the same place.

\medskip
\noindent\fbox{\parbox{0.97\textwidth}{\textbf{Proposed rule.} Two candidate records (same
location rounded to 4 decimals and same \col{main\_value}) are considered ``very probably the same real place''
when
\[
\text{name\_sim} \;\geq\; 0.66 ,
\]
i.e.\ at least about two-thirds of the name words are shared. \col{tag\_sim} is \emph{not}
used, because it is uncorrelated with real identity and would reject the true node + way
duplicates (whose \col{tag\_sim} is 0).}}

\medskip
\noindent\textbf{Result: 532 of the 2{,}396 pairs} are marked as duplicates, involving
\textbf{1{,}014 records}. Because words are compared as sets, a change in word order is also
detected (\fa{آنلاین موبایل} $\equiv$ \fa{موبایل آنلاین}).
\emph{Limitation:} very generic names such as \fa{آلاچیق} (shelter) give
$\text{name\_sim}=1$ even for two different shelters a few metres apart. A stricter rule
could also require a minimum name length or a distance in metres.

% =====================================================================
\section{Step 6: Descriptive Statistics, Visualization and OLAP}

\subsection{Frequency and mode of \col{main\_value} within each \col{main\_key}}

\begin{table}[H]
\centering\footnotesize
\setlength{\tabcolsep}{3.5pt}
\caption{Mode of \col{main\_value} in each category.}
\begin{tabular}{lrrlrrl}
\toprule
\col{main\_key} & \textbf{Records} & \textbf{Distinct} & \textbf{Mode} & \textbf{Count} & \textbf{\%} & \textbf{Top 3} \\
\midrule
amenity    & 138{,}493 & 257 & place\_of\_worship & 17{,}810 & 12.9 & place\_of\_worship, school, bank \\
shop       & 72{,}122  & 543 & supermarket        & 7{,}667  & 10.6 & supermarket, bakery, car\_repair \\
leisure    & 42{,}320  & 56  & park               & 20{,}097 & 47.5 & park, pitch, garden \\
office     & 25{,}972  & 181 & government         & 9{,}389  & 36.2 & government, company, estate\_agent \\
tourism    & 12{,}220  & 30  & hotel              & 2{,}577  & 21.1 & hotel, attraction, information \\
craft      & 3{,}911   & 107 & photographer       & 366      & 9.4  & photographer, electronics\_repair, carpenter \\
healthcare & 2{,}114   & 34  & laboratory         & 957      & 45.3 & laboratory, yes, physiotherapist \\
\bottomrule
\end{tabular}
\end{table}

\begin{table}[H]
\centering
\caption{The 10 most frequent \col{amenity} types.}
\begin{tabular}{lr@{\hspace{1.2cm}}lr}
\toprule
place\_of\_worship & 17{,}810 & fuel     & 5{,}721 \\
school             & 16{,}267 & clinic   & 4{,}692 \\
bank               & 11{,}943 & shelter  & 4{,}646 \\
parking            & 9{,}699  & cafe     & 4{,}517 \\
restaurant         & 9{,}006  & pharmacy & 4{,}425 \\
\bottomrule
\end{tabular}
\end{table}

\begin{itemize}
  \item The most frequent type of \textbf{amenity} is \textbf{place\_of\_worship} (17{,}810, 12.9\%),
  followed by school and bank.
  \item \col{leisure} and \col{healthcare} are dominated by their mode: park makes up almost
  half of leisure, and laboratory 45\% of healthcare.
  \item \col{shop} (543 distinct values) and \col{amenity} (257) are very diverse.
  \item The value \col{yes} among the top healthcare types (\col{healthcare=yes}) is another
  quality issue: the mapper did not say what kind of healthcare place it is.
\end{itemize}

\subsection{Percentiles of \col{hub\_km}}

\begin{table}[H]
\centering
\caption{Percentiles of \col{hub\_km}.}
\begin{tabular}{cccc}
\toprule
$P_{25}$ & $P_{50}$ & $P_{75}$ & $P_{90}$ \\
\midrule
5.00 km & 14.99 km & 69.98 km & 120.95 km \\
\bottomrule
\end{tabular}
\end{table}

Half of all POIs lie within 15\,km of an approximate provincial center. This shows that OSM
coverage is concentrated in and around the big cities. The distribution is strongly
right-skewed (mean 45.3\,km is much larger than the median of 15\,km).

\textbf{Records farther than 300\,km: 2{,}300 (0.77\%).}

\begin{table}[H]
\centering
\caption{Records with \col{hub\_km} $>300$.}
\begin{tabular}{lr@{\hspace{1.5cm}}lr}
\toprule
\col{nearest\_hub} & \textbf{Count} & \col{addr:city} (top) & \textbf{Count} \\
\midrule
Zahedan     & 1{,}784 & \fa{نیک شهر}   & 49 \\
BandarAbbas & 515     & \fa{چابهار}    & 20 \\
Ardabil     & 1       & \fa{شهر چابهار} & 10 \\
            &         & \fa{مخت}       & 6 \\
            &         & \fa{قصرقند}    & 3 \\
\bottomrule
\end{tabular}
\end{table}

\begin{figure}[H]
\centering
\includegraphics[width=0.55\textwidth]{figures/far_map.png}
\caption{All records (lon/lat), with the records more than 300\,km from their nearest hub in red.}
\end{figure}

\textbf{Where are they?} The far-away records have a median latitude of 25.99 and a median
longitude of 60.63 (interquartile ranges 25.35--26.38 and 60.20--61.14). This is
\textbf{southern Sistan and Baluchestan} along the Makran coast (Chabahar, Nikshahr, Konarak,
Qasr-e Qand). This area is very far from every hub in the list: Zahedan is in the far north
of that large province, and Bandar Abbas is in the next province to the west. The result
comes from using straight-line distance to a few approximate centers instead of real
administrative boundaries. One record (nearest hub Ardabil) has latitude $\approx 43$,
which is outside Iran's borders: either a coordinate error or a point outside the country in
the extract.

\subsection{Plots}

\begin{figure}[H]
\centering
\includegraphics[width=0.78\textwidth]{figures/year_hist.png}
\caption{Histogram of the year of \col{edit\_date}.}
\end{figure}

\begin{table}[H]
\centering\small
\caption{Number of records per edit year.}
\begin{tabular}{*{10}{c}}
\toprule
2007 & 2008 & 2009 & 2010 & 2011 & 2012 & 2013 & 2014 & 2015 & 2016 \\
3 & 6 & 11 & 221 & 376 & 389 & 661 & 1{,}580 & 2{,}835 & 11{,}966 \\
\midrule
2017 & 2018 & 2019 & 2020 & 2021 & 2022 & 2023 & 2024 & 2025 & 2026 \\
21{,}156 & 21{,}921 & 32{,}808 & 18{,}109 & 24{,}011 & 26{,}366 & 44{,}647 & 38{,}500 & 35{,}096 & 16{,}490 \\
\bottomrule
\end{tabular}
\end{table}

\textbf{Interpretation.} OSM coverage of Iran grew in stages:
\begin{itemize}
  \item there are almost no records before 2010;
  \item growth is slow until 2015;
  \item there is a jump in \textbf{2016} (about 12k), then 20k--45k records per year, with a peak in \textbf{2023} (44{,}647);
  \item the dip in 2020 matches the COVID year;
  \item 2026 is lower because the year was not finished when the snapshot was taken.
\end{itemize}
\emph{Caveat:} \col{edit\_date} is the date of the \emph{last} edit, not of creation. A place
created in 2012 and edited in 2023 is counted in 2023, so the histogram is shifted toward
recent years and shows mapping \emph{activity} rather than pure creation.

\begin{figure}[H]
\centering
\includegraphics[width=0.75\textwidth]{figures/hub_boxplot.png}
\caption{Boxplot of \col{hub\_km} for the five most frequent \col{nearest\_hub} values.}
\end{figure}

\begin{table}[H]
\centering
\caption{\col{hub\_km} statistics for the five most frequent hubs.}
\begin{tabular}{lrrrrrrr}
\toprule
\col{nearest\_hub} & \textbf{Count} & \textbf{Mean} & \textbf{Std} & \textbf{25\%} & \textbf{50\%} & \textbf{75\%} & \textbf{Max} \\
\midrule
Tehran  & 60{,}730 & 13.5 & 14.0 & 5.8 & 9.8  & 13.9 & 119.8 \\
Isfahan & 21{,}433 & 23.4 & 29.8 & 3.6 & 9.4  & 30.4 & 193.4 \\
Shiraz  & 21{,}406 & 47.0 & 65.8 & 6.1 & 10.7 & 60.4 & 285.4 \\
Karaj   & 17{,}876 & 30.4 & 31.2 & 6.9 & 18.7 & 40.8 & 104.8 \\
Mashhad & 16{,}375 & 38.8 & 53.4 & 4.1 & 9.3  & 71.4 & 263.4 \\
\bottomrule
\end{tabular}
\end{table}

\textbf{Interpretation.}
\begin{itemize}
  \item \textbf{Tehran} has the most records and the most compact distribution (median 9.8\,km,
  IQR 5.8--13.9\,km), because the region is small and densely mapped.
  \item \textbf{Shiraz} and \textbf{Mashhad} have a similar median (about 10\,km) but much longer
  boxes and whiskers ($P_{75}\approx 60$--$71$\,km, maximum 263--285\,km), because they serve
  large provinces with many distant towns.
  \item \textbf{Karaj} has the highest median (18.7\,km), because it takes in the many suburbs
  between Tehran and Qazvin.
  \item All five distributions are right-skewed, with many outliers above the upper whisker.
\end{itemize}

\subsection{Data cube and OLAP operations}

A data cube with three dimensions, \textbf{\col{nearest\_hub} $\times$ \col{main\_key}
$\times$ year of \col{edit\_date}}, was built. Each cell holds the \textbf{number of records}.
It is stored as a pandas Series with a three-level MultiIndex. Only non-empty cells are kept,
which suits a sparse cube:

\begin{lstlisting}
df["year"] = pd.to_datetime(df["edit_date"]).dt.year
cube = df.groupby(["nearest_hub", "main_key", "year"]).size()
\end{lstlisting}

Only \textbf{2{,}722 of the 4{,}340} possible cells (31 hubs $\times$ 7 categories $\times$ 20
years) are non-empty, so the cube is about 37\% empty, mainly because of the early years and
the rare categories in small hubs.

\paragraph{Roll-up (remove the year dimension).}
\begin{lstlisting}
rollup = cube.groupby(level=["nearest_hub", "main_key"]).sum().unstack(fill_value=0)
\end{lstlisting}

\begin{table}[H]
\centering\small
\caption{Roll-up: records per hub and category over all years (top 10 hubs).}
\begin{tabular}{lrrrrrrrr}
\toprule
\col{nearest\_hub} & amenity & craft & healthcare & leisure & office & shop & tourism & \textbf{total} \\
\midrule
Tehran      & 26{,}145 & 792 & 778 & 9{,}802 & 6{,}156 & 15{,}768 & 1{,}289 & \textbf{60{,}730} \\
Isfahan     & 9{,}562  & 387 & 128 & 2{,}758 & 1{,}677 & 5{,}818  & 1{,}103 & \textbf{21{,}433} \\
Shiraz      & 8{,}607  & 242 & 134 & 2{,}731 & 1{,}807 & 7{,}210  & 675     & \textbf{21{,}406} \\
Karaj       & 7{,}215  & 435 & 111 & 2{,}994 & 1{,}817 & 4{,}834  & 470     & \textbf{17{,}876} \\
Mashhad     & 7{,}257  & 209 & 113 & 1{,}901 & 1{,}289 & 4{,}445  & 1{,}161 & \textbf{16{,}375} \\
Tabriz      & 6{,}505  & 172 & 76  & 1{,}329 & 1{,}111 & 2{,}712  & 482     & \textbf{12{,}387} \\
Sari        & 4{,}593  & 166 & 50  & 924     & 886     & 3{,}282  & 364     & \textbf{10{,}265} \\
Rasht       & 4{,}257  & 123 & 61  & 671     & 853     & 3{,}013  & 594     & \textbf{9{,}572} \\
BandarAbbas & 4{,}783  & 120 & 41  & 1{,}297 & 707     & 2{,}074  & 550     & \textbf{9{,}572} \\
Arak        & 4{,}616  & 138 & 25  & 1{,}310 & 731     & 2{,}280  & 422     & \textbf{9{,}522} \\
\bottomrule
\end{tabular}
\end{table}

\noindent\textit{Question answered:} ``Over all years, how many POIs of each category does each
hub region have?'' The Tehran region is far ahead (60{,}730, about three times Isfahan or Shiraz),
and \col{amenity} is the largest category in every hub.

\paragraph{Drill-down (Tehran / \col{leisure} $\rightarrow$ \col{main\_value}).}
A finer cube that also has the \col{main\_value} level under \col{main\_key} was built, and
the cell (Tehran, leisure) was opened:

\begin{lstlisting}
cube_detail = df.groupby(["nearest_hub", "main_key", "main_value", "year"]).size()
tehran_leisure = cube_detail.loc[("Tehran", "leisure")].groupby(level="main_value").sum()
\end{lstlisting}

\begin{table}[H]
\centering
\caption{Drill-down: sub-types of the 9{,}802 \col{leisure} records in Tehran (top 10).}
\begin{tabular}{lr@{\hspace{1.2cm}}lr}
\toprule
park          & 3{,}490 & playground       & 527 \\
garden        & 2{,}252 & fitness\_centre  & 254 \\
pitch         & 1{,}745 & stadium          & 70 \\
swimming\_pool & 703    & fitness\_station & 43 \\
sports\_centre & 538    & amusement\_arcade & 41 \\
\bottomrule
\end{tabular}
\end{table}

\noindent\textit{Question answered:} ``What does leisure in Tehran actually consist of?''
It is mainly parks, gardens and sports pitches (together 76\%), followed by swimming pools,
sports centres and playgrounds.

\paragraph{Selection (\col{nearest\_hub} = Tehran).}
\begin{lstlisting}
tehran_slice = cube.xs("Tehran", level="nearest_hub")
tehran_dist = tehran_slice.groupby(level="main_key").sum()
\end{lstlisting}

\begin{table}[H]
\centering
\caption{Selection: category distribution inside the Tehran region compared with all of Iran.}
\begin{tabular}{lrrr}
\toprule
\col{main\_key} & \textbf{Count} & \textbf{Tehran \%} & \textbf{Whole Iran \%} \\
\midrule
amenity    & 26{,}145 & 43.1 & 46.61 \\
shop       & 15{,}768 & 26.0 & 24.27 \\
leisure    & 9{,}802  & 16.1 & 14.24 \\
office     & 6{,}156  & 10.1 & 8.74 \\
tourism    & 1{,}289  & 2.1  & 4.11 \\
craft      & 792      & 1.3  & 1.32 \\
healthcare & 778      & 1.3  & 0.71 \\
\bottomrule
\end{tabular}
\end{table}

\noindent\textit{Question answered:} ``How are the POI categories distributed in the Tehran region,
and how does that compare with the whole country?'' \col{amenity} is still the largest share
(43.1\%). Compared with all of Iran, Tehran has relatively more \col{shop}, \col{leisure},
\col{office} and especially \col{healthcare} (1.3\% vs 0.71\%), and about half the share of
\col{tourism} (2.1\% vs 4.11\%). This fits a capital that is a business and medical center
rather than a tourist destination.

% =====================================================================
\section{Problems Encountered and Solutions}

\begin{table}[H]
\centering
\begin{tabularx}{\textwidth}{>{\raggedright\arraybackslash}X>{\raggedright\arraybackslash}X}
\toprule
\textbf{Problem} & \textbf{Solution} \\
\midrule
pandas would auto-convert \col{phone}/\col{capacity} and lose leading \col{0} and \col{+}.
& Read these columns with \col{dtype=str}. \\
Persian (\fa{۰}--\fa{۹}) and Arabic-Indic digits in numeric fields.
& One \col{str.maketrans} table that maps both to Latin digits. \\
Arabic \fa{ي}/\fa{ك} versus Persian \fa{ی}/\fa{ک}.
& Replaced before comparing cities and names. \\
A city written in many forms (\col{Tehran}, \fa{شهر تهران}, \dots).
& Whitespace collapsing, lower-casing, removing the \fa{شهر} prefix, Latin$\rightarrow$Persian
dictionary (2{,}525 $\rightarrow$ 2{,}281). \\
Substring search for Tehran also matched other places (\fa{تهرانسر}).
& For the normalization, only exact rules and a dictionary were used, not substring matching. \\
Several phone numbers in one field.
& Split on \col{;} \col{,} \col{/} and check every part. \\
Different country-code / trunk-prefix conventions.
& \col{nationals()} reduces every number to its national significant digits before the length check. \\
The self-join produced self-pairs and mirrored pairs.
& Keep only rows with \col{index\_1 < index\_2}. \\
Jaccard of two empty tag sets (0/0).
& Return \col{NaN} instead of dividing by zero. \\
Tag-key similarity turned out to be uninformative for duplicates.
& The duplicate threshold uses name-word similarity only. \\
Mixed integers and missing values in \col{capacity}.
& The nullable integer dtype \col{Int64}. \\
\bottomrule
\end{tabularx}
\end{table}

% =====================================================================
\section{Use of Large Language Models}

An LLM (Claude) was used to add explanatory Markdown cells and code
comments to the notebook, and to draft the text and layout of this \LaTeX{} report from the
notebook's code and outputs. All numbers in the report come directly from the
outputs of the notebook.

% =====================================================================
\section{Conclusion}

The audit confirms that \col{iran\_poi\_raw.csv} has the typical problems of volunteered
geographic information:
\begin{itemize}
  \item free-text fields with many spellings for the same thing (2{,}525 city values);
  \item mixed Persian/Arabic characters and digits;
  \item numeric fields that contain text, units, negative values and placeholders;
  \item phone numbers in at least nine formats;
  \item thousands of near-duplicate records, mostly from mapping a place both as a point and as a building.
\end{itemize}
Simple, explicit rules were enough to fix a large part of this:
\begin{itemize}
  \item character and digit unification and a city dictionary;
  \item a strict integer parser for \col{capacity};
  \item a length rule for phones;
  \item a hand-written Jaccard similarity on name words, which flagged 532 duplicate pairs.
\end{itemize}
The descriptive and OLAP analysis shows a dataset concentrated around a few big cities
(especially Tehran) and growing quickly since 2016. Any regional conclusion depends on the
approximate, distance-based \col{nearest\_hub} assignment and not on official boundaries.

\end{document}
